\documentclass{article}
\usepackage{amsmath,amssymb}
\usepackage[margin=1in]{geometry}
\usepackage{enumitem}

% -------------------------------------------------
% Color + framed boxes (from lecture notes)
% -------------------------------------------------
\usepackage{xcolor}
\definecolor{myblue}{RGB}{0,80,180}

\newcommand{\bluebox}[1]{%
  \noindent\begingroup
  \setlength{\fboxsep}{8pt}%
  \setlength{\fboxrule}{1.2pt}%
  \fcolorbox{myblue}{white}{%
    \parbox{\dimexpr\textwidth-2\fboxsep-2\fboxrule\relax}{\color{black}#1}%
  }%
  \endgroup
}

\begin{document}


\begin{flushright}
Neil Hwang\\
Linear Algebra\\
March 25, 2026\\
Homework 1\\

\end{flushright}




\bigskip

\hfill \break  $~\!\!$ \dotfill \medskip \\
\textbf{Question.}\\
Let $W \subset V$. Show that $W$ is a subspace of $V$ if and only if:
\begin{align*}
\forall x,y \in W,\quad x+y \in W, \\
\forall a \in F,\ \forall x \in W,\quad ax \in W, \\
0 \in W.
\end{align*}

\noindent\\
\textbf{Solution.}\\
\quad\\
\noindent\fbox{ \parbox{\textwidth}{
Suppose first that $W$ is a subspace of $V$. Then $W$ is itself a vector space with the same operations as $V$. Hence:
\begin{align*}
x,y\in W &\implies x+y\in W, \\
a\in F,\ x\in W &\implies ax\in W, \\
0&\in W.
\end{align*}

Conversely, assume that $W\subset V$ satisfies
\begin{align*}
x,y\in W &\implies x+y\in W, \\
a\in F,\ x\in W &\implies ax\in W, \\
0&\in W.
\end{align*}
We show that $W$ is a subspace of $V$.

Since $W\subset V$, the addition and scalar multiplication on $W$ are the same as those on $V$. Because $V$ is a vector space, the associative, commutative, and distributive laws automatically hold on $W$ as well.

Also, for any $x\in W$, since $-1\in F$ and $W$ is closed under scalar multiplication,
\begin{align*}
-x = (-1)x \in W.
\end{align*}
Thus every element of $W$ has an additive inverse in $W$.

The zero vector belongs to $W$ by assumption. Closure under addition is already assumed, and closure under scalar multiplication is also assumed. Therefore all vector space axioms hold on $W$, so $W$ is a subspace of $V$.
}}

\medskip

\bluebox{
\textbf{Alternative solution.} Note that the condition $0 \in W$ is redundant when closure under scalar multiplication is assumed. If $W \neq \varnothing$, pick any $w \in W$. Then $0 = 0 \cdot w \in W$ by scalar multiplication closure. So the subspace test from the lecture notes uses only two conditions:
\begin{align*}
\text{(i)}\; x,y \in W &\Rightarrow x+y \in W, \\
\text{(ii)}\; a \in F,\, x \in W &\Rightarrow ax \in W,
\end{align*}
together with $W \neq \varnothing$.

Given these, the zero vector, additive inverses ($-x = (-1)x$), and all remaining axioms are inherited from $V$. Hence $W$ is a subspace.
}

\bigskip


\hfill \break  $~\!\!$ \dotfill \medskip \\
\textbf{Question.}\\
Let $T: V \to W$ be a function. Show that $T$ is linear if and only if:
\begin{align*}
T(x+y) = T(x) + T(y), \\
T(ax) = aT(x)
\end{align*}
for all $x,y \in V$ and $a \in F$.

\noindent\\
\textbf{Solution.}\\
\quad\\
\noindent\fbox{ \parbox{\textwidth}{
By definition, $T$ is linear if for all $x,y\in V$ and all $a,b\in F$,
\begin{align*}
T(ax+by)=aT(x)+bT(y).
\end{align*}

Suppose first that $T$ is linear. Then taking $a=b=1$, we get
\begin{align*}
T(x+y)=T(1x+1y)=1T(x)+1T(y)=T(x)+T(y).
\end{align*}
Taking $b=0$ and $y=0$, we get
\begin{align*}
T(ax)=T(ax+0\cdot 0)=aT(x)+0\cdot T(0)=aT(x).
\end{align*}
So the two displayed properties follow.

Conversely, suppose that for all $x,y\in V$ and all $a\in F$,
\begin{align*}
T(x+y)=T(x)+T(y),\qquad T(ax)=aT(x).
\end{align*}
We show that $T$ is linear. Let $x,y\in V$ and $a,b\in F$. Then
\begin{align*}
T(ax+by)
&=T(ax)+T(by) \\
&=aT(x)+bT(y).
\end{align*}
Thus $T$ satisfies the defining property of linearity, hence $T$ is linear.
}}

\medskip

\bluebox{
\textbf{Alternative solution.} One can combine both conditions into a single check. $T$ is linear if and only if
\[
T(\alpha x + \beta y) = \alpha T(x) + \beta T(y)
\]
for all $x, y \in V$ and $\alpha, \beta \in F$. To see the equivalence directly:

$(\Rightarrow)$: Apply the single condition with $\alpha = \beta = 1$ to get additivity, and with $\beta = 0$ to get homogeneity.

$(\Leftarrow)$: Given additivity and homogeneity, compute:
\[
T(\alpha x + \beta y) = T(\alpha x) + T(\beta y) = \alpha T(x) + \beta T(y).
\]
This shows the two-condition formulation and the one-condition formulation are logically equivalent. In practice, verifying the two conditions separately is often easier.
}

\bigskip

\hfill \break  $~\!\!$ \dotfill \medskip \\
\textbf{Question.}\\
Let $B = \{x_1,\dots,x_n\}$ be a basis of $V$. Show that every $x \in V$ can be written uniquely as:
\begin{align*}
x = a_1 x_1 + \cdots + a_n x_n.
\end{align*}

\noindent\\
\textbf{Solution.}\\
\quad\\
\noindent\fbox{ \parbox{\textwidth}{
Since $B$ is a basis of $V$, by definition it spans $V$ and is linearly independent. Because $B$ spans $V$, every $x\in V$ can be expressed as a linear combination of the basis vectors:
\begin{align*}
x=a_1x_1+\cdots+a_nx_n
\end{align*}
for some scalars $a_1,\dots,a_n\in F$. This proves existence.

To prove uniqueness, suppose that
\begin{align*}
x=a_1x_1+\cdots+a_nx_n=b_1x_1+\cdots+b_nx_n.
\end{align*}
Subtracting the two expressions gives
\begin{align*}
0=(a_1-b_1)x_1+\cdots+(a_n-b_n)x_n.
\end{align*}
Since $B$ is linearly independent, the only linear combination of $x_1,\dots,x_n$ equal to $0$ is the trivial one. Therefore
\begin{align*}
a_1-b_1=0,\quad \dots,\quad a_n-b_n=0.
\end{align*}
Hence
\begin{align*}
a_i=b_i\qquad \text{for all } i=1,\dots,n.
\end{align*}
So the representation is unique.
}}

\medskip

\bluebox{
\textbf{Alternative solution.} Use the coordinate map. Define $\Phi_B : V \to F^n$ by $\Phi_B(x) = [x]_B$. Since $B$ is a basis, $\Phi_B$ is a linear isomorphism (as shown in the lecture notes). In particular, $\Phi_B$ is a bijection. But a bijection is injective, and injectivity of $\Phi_B$ means: if $\Phi_B(x) = (a_1, \dots, a_n)$ and also $\Phi_B(x) = (b_1, \dots, b_n)$, then the two tuples must be identical. This is exactly uniqueness of the representation $x = a_1 x_1 + \cdots + a_n x_n$.

Surjectivity gives existence: every tuple $(a_1, \dots, a_n) \in F^n$ arises as the coordinate vector of some $x \in V$, so every $x$ has at least one representation.
}

\bigskip

\hfill \break  $~\!\!$ \dotfill \medskip \\
\textbf{Question.}\\
Let $V$ be an $n$-dimensional vector space with ordered basis $B = \{x_1,\dots,x_n\}$. Define the coordinate map:
\begin{align*}
[x]_B = (a_1,\dots,a_n) \in F^n
\end{align*}
where $x = a_1 x_1 + \cdots + a_n x_n$.

Show that the map $x \mapsto [x]_B$ is a linear isomorphism from $V$ to $F^n$.

\noindent\\
\textbf{Solution.}\\
\quad\\
\noindent\fbox{ \parbox{\textwidth}{
Define $\Phi:V\to F^n$ by
\begin{align*}
\Phi(x)=[x]_B,
\end{align*}
where if $x=a_1x_1+\cdots+a_nx_n,$ then $[x]_B=(a_1,\dots,a_n).$

We first show that $\Phi$ is linear. Let
\begin{align*}
x=a_1x_1+\cdots+a_nx_n,\qquad y=b_1x_1+\cdots+b_nx_n.
\end{align*}
Then
\begin{align*}
x+y=(a_1+b_1)x_1+\cdots+(a_n+b_n)x_n,
\end{align*}
so
\begin{align*}
[x+y]_B=(a_1+b_1,\dots,a_n+b_n)=[x]_B+[y]_B.
\end{align*}
Also, for $c\in F$,
\begin{align*}
cx=(ca_1)x_1+\cdots+(ca_n)x_n,
\end{align*}
hence
\begin{align*}
[cx]_B=(ca_1,\dots,ca_n)=c[x]_B.
\end{align*}
Therefore $\Phi$ is linear.

Next we show injectivity. Suppose $\Phi(x)=\Phi(y)$. Then
\begin{align*}
[x]_B=[y]_B.
\end{align*}
This means the coordinates of $x$ and $y$ with respect to $B$ are identical. Since basis expansions are unique, it follows that $x=y$. Hence $\Phi$ is injective.

Now we show surjectivity. Let $(a_1,\dots,a_n)\in F^n$. Define
\begin{align*}
x=a_1x_1+\cdots+a_nx_n\in V.
\end{align*}
Then by definition,
\begin{align*}
\Phi(x)=[x]_B=(a_1,\dots,a_n).
\end{align*}
Thus every vector in $F^n$ is in the image of $\Phi$, so $\Phi$ is surjective.
Note that in linear algebra, ``linear'' = ``homomorphism.'' It is exactly what it means for $\Phi$ to be a homomorphism of vector spaces.

Since $\Phi$ is linear, injective, and surjective, it is a linear isomorphism.
}}

\medskip

\bluebox{
\textbf{Alternative solution.} A dimension-counting argument shortens the proof. Since $\Phi_B$ is linear (verified as above), it suffices to show $\ker(\Phi_B) = \{0\}$ and then invoke the dimension theorem.

If $\Phi_B(x) = 0$, then $[x]_B = (0, \dots, 0)$, meaning $x = 0 \cdot x_1 + \cdots + 0 \cdot x_n = 0$. So $\ker(\Phi_B) = \{0\}$, and $\Phi_B$ is injective.

By the rank--nullity theorem:
\[
\dim(V) = \dim(\ker \Phi_B) + \dim(\operatorname{Im} \Phi_B) \implies n = 0 + \dim(\operatorname{Im} \Phi_B).
\]
So $\dim(\operatorname{Im} \Phi_B) = n = \dim(F^n)$, hence $\operatorname{Im} \Phi_B = F^n$, giving surjectivity. Therefore $\Phi_B$ is a linear isomorphism.
}

\bigskip

\hfill \break  $~\!\!$ \dotfill \medskip \\
\textbf{Question.}\\
Let $B, B'$ be two ordered bases of $V$. Show that there exists an invertible matrix $Q$ such that:
\begin{align*}
[x]_{B'} = Q [x]_B
\end{align*}
and identify $Q$ as the change of coordinate matrix.

\noindent\\
\textbf{Solution.}\\
\quad\\
\noindent\fbox{ \parbox{\textwidth}{
Let
\begin{align*}
B=\{x_1,\dots,x_n\},\qquad B'=\{x_1',\dots,x_n'\}.
\end{align*}
Consider the two coordinate isomorphisms
\begin{align*}
\Phi_B:V\to F^n,\qquad \Phi_B(x)=[x]_B,
\end{align*}
and
\begin{align*}
\Phi_{B'}:V\to F^n,\qquad \Phi_{B'}(x)=[x]_{B'}.
\end{align*}
Both are linear isomorphisms. Therefore the composition
\begin{align*}
Q:=\Phi_{B'}\circ \Phi_B^{-1}:F^n\to F^n
\end{align*}
is a linear isomorphism. Hence $Q$ is represented by an invertible $n\times n$ matrix.

Now let $x\in V$. Then
\begin{align*}
Q[x]_B
&=(\Phi_{B'}\circ \Phi_B^{-1})([x]_B) \\
&=\Phi_{B'}(x) \\
&=[x]_{B'}.
\end{align*}
So
\begin{align*}
[x]_{B'}=Q[x]_B.
\end{align*}

Here, this matrix $Q$ is identified as the change of coordinate matrix from $B$-coordinates to $B'$-coordinates.

Equivalently, the $j$th column of $Q$ is the coordinate vector of the basis vector $x_j$ with respect to $B'$, namely
\begin{align*}
Q=\begin{bmatrix}
[x_1]_{B'} & [x_2]_{B'} & \cdots & [x_n]_{B'}
\end{bmatrix}.
\end{align*}
Since $Q$ comes from a linear isomorphism, it is invertible.
}}

\medskip

\bluebox{
\textbf{Alternative solution.} Construct $Q$ directly by column expansion. Write each old basis vector $x_j \in B$ in terms of $B' = \{x_1', \dots, x_n'\}$:
\[
x_j = q_{1j} x_1' + \cdots + q_{nj} x_n', \quad j = 1, \dots, n.
\]
Then $[x_j]_{B'} = (q_{1j}, \dots, q_{nj})^T$. Set $Q$ to be the matrix whose $j$th column is $[x_j]_{B'}$.

For any $x = a_1 x_1 + \cdots + a_n x_n$, substituting gives
\[
x = \sum_{j=1}^n a_j \left( \sum_{i=1}^n q_{ij} x_i' \right) = \sum_{i=1}^n \left( \sum_{j=1}^n q_{ij} a_j \right) x_i'.
\]
Reading off coordinates in $B'$: $[x]_{B'} = Q [x]_B$.

Invertibility: if $Q c = 0$, then $c_1 x_1 + \cdots + c_n x_n = 0$ (since the coordinate map is injective), and linear independence of $B$ forces $c = 0$.
}

\bigskip

\hfill \break  $~\!\!$ \dotfill \medskip \\
\textbf{Question.}\\
Let $T: V \to V$ be linear and let $B, B'$ be bases. Show that:
\begin{align*}
[T]_{B'} = Q [T]_B Q^{-1}
\end{align*}
where $Q$ is the change of coordinate matrix from $B'$ to $B$.

\noindent\\
\textbf{Solution.}\\
\quad\\
\noindent\fbox{ \parbox{\textwidth}{
Let
\begin{align*}
\Phi_B(x)=[x]_B,\qquad \Phi_{B'}(x)=[x]_{B'}.
\end{align*}
Then the matrix of $T$ with respect to $B$ is the unique matrix $[T]_B$ satisfying
\begin{align*}
[T(x)]_B=[T]_B[x]_B\qquad \text{for all }x\in V.
\end{align*}
Similarly,
\begin{align*}
[T(x)]_{B'}=[T]_{B'}[x]_{B'}.
\end{align*}

Let $Q$ be the change of coordinate matrix from $B'$ to $B$, so
\begin{align*}
[x]_B=Q[x]_{B'}.
\end{align*}
Hence
\begin{align*}
[x]_{B'}=Q^{-1}[x]_B.
\end{align*}

Now for any $x\in V$,
\begin{align*}
[T(x)]_{B'}
&=Q^{-1}[T(x)]_B \\
&=Q^{-1}[T]_B[x]_B \\
&=Q^{-1}[T]_BQ[x]_{B'}.
\end{align*}
Since also
\begin{align*}
[T(x)]_{B'}=[T]_{B'}[x]_{B'},
\end{align*}
we conclude that
\begin{align*}
[T]_{B'}=Q^{-1}[T]_BQ.
\end{align*}

If instead one defines $Q$ to be the change of coordinate matrix from $B$ to $B'$, then the formula becomes
\begin{align*}
[T]_{B'}=Q[T]_BQ^{-1}.
\end{align*}
where the exact placement of $Q$ and $Q^{-1}$ depends on the convention for the change of coordinate matrix.
}}

\medskip

\bluebox{
\textbf{Alternative solution (commutative diagram).} From the lecture notes, the key idea is that the following diagram commutes:
\[
\begin{array}{ccc}
V & \xrightarrow{T} & V \\
\downarrow [\cdot]_{B'} & & \downarrow [\cdot]_{B'} \\
F^n & \xrightarrow{[T]_{B'}} & F^n
\end{array}
\]
Tracing both routes from top-left to bottom-right:

\textbf{Route 1 (top then right):} $x \xrightarrow{T} T(x) \xrightarrow{[\cdot]_{B'}} [T(x)]_{B'}$.

\textbf{Route 2 (left then bottom):} $x \xrightarrow{[\cdot]_{B'}} [x]_{B'} \xrightarrow{[T]_{B'}} [T]_{B'} [x]_{B'}$.

Now $[\cdot]_{B'} = Q^{-1} \circ [\cdot]_B$ (since $[x]_B = Q[x]_{B'}$ implies $[x]_{B'} = Q^{-1}[x]_B$). Substituting into Route 1:
\[
[T(x)]_{B'} = Q^{-1}[T(x)]_B = Q^{-1} [T]_B [x]_B = Q^{-1} [T]_B Q [x]_{B'}.
\]
Comparing with Route 2 gives $[T]_{B'} = Q^{-1}[T]_B Q$. This is a \emph{similarity transformation}: the same linear operator, viewed through different coordinate systems.
}

\bigskip

\hfill \break  $~\!\!$ \dotfill \medskip \\
\textbf{Question.}\\
Fill in the missing steps to prove that if $T,S$ are linear, then $T+S$ is linear:
\begin{align*}
(T+S)(x+y) &= (T+S)(x) + (T+S)(y), \\
(T+S)(ax) &= a(T+S)(x).
\end{align*}

\noindent\\
\textbf{Solution.}\\
\quad\\
\noindent\fbox{ \parbox{\textwidth}{
Let $x,y\in X$ and $a\in F$. Then
\begin{align*}
(T+S)(x+y)
&=T(x+y)+S(x+y) \\
&=(T(x)+T(y))+(S(x)+S(y)) \\
&=(T(x)+S(x))+(T(y)+S(y)) \\
&=(T+S)(x)+(T+S)(y).
\end{align*}
Also,
\begin{align*}
(T+S)(ax)
&=T(ax)+S(ax) \\
&=aT(x)+aS(x) \\
&=a(T(x)+S(x)) \\
&=a(T+S)(x).
\end{align*}
Hence $T+S$ is linear.
}}

\medskip

\bluebox{
\textbf{Alternative solution.} Use the combined linearity condition. It suffices to show that for all $x, y \in X$ and $\alpha, \beta \in F$:
\[
(T+S)(\alpha x + \beta y) = \alpha (T+S)(x) + \beta (T+S)(y).
\]
Compute:
\begin{align*}
(T+S)(\alpha x + \beta y)
&= T(\alpha x + \beta y) + S(\alpha x + \beta y) \\
&= [\alpha T(x) + \beta T(y)] + [\alpha S(x) + \beta S(y)] \\
&= \alpha [T(x) + S(x)] + \beta [T(y) + S(y)] \\
&= \alpha (T+S)(x) + \beta (T+S)(y).
\end{align*}
Each step uses linearity of $T$ and $S$ individually, then rearranges using vector space axioms in $Y$.
}

\bigskip

\hfill \break  $~\!\!$ \dotfill \medskip \\
\textbf{Question.}\\
Show that checking closure under addition and scalar multiplication is sufficient to prove $W \subset V$ is a subspace. Explain why $0 \in W$ automatically follows.

\quad\\
\noindent
\textbf{Solution.}\\
\quad\\
\noindent\fbox{ \parbox{\textwidth}{
Assume $W\neq \varnothing$, and assume:
\begin{align*}
x,y\in W &\implies x+y\in W, \\
a\in F,\ x\in W &\implies ax\in W.
\end{align*}
We show that $W$ is a subspace.

Pick any $w\in W$; this is possible because $W\neq\varnothing$. Since $0\in F$, closure under scalar multiplication gives
\begin{align*}
0\cdot w=0\in W.
\end{align*}
Thus the zero vector belongs to $W$ automatically.

Next, if $w\in W$, then since $-1\in F$,
\begin{align*}
(-1)w=-w\in W.
\end{align*}
So additive inverses are also in $W$.

Because all remaining vector space axioms are inherited from $V$, it follows that $W$ is a subspace of $V$.

Therefore, for a nonempty subset $W$, closure under addition and scalar multiplication is enough to conclude that $W$ is a subspace.
}}

\medskip

\bluebox{
\textbf{Alternative solution.} One can combine the two closure conditions into a single condition. $W$ is a subspace if and only if $W \neq \varnothing$ and
\[
\alpha x + \beta y \in W \quad \text{for all } x, y \in W \text{ and } \alpha, \beta \in F.
\]
To see that $0 \in W$: pick any $w \in W$ and take $\alpha = \beta = 0$, $x = y = w$. Then $0 \cdot w + 0 \cdot w = 0 \in W$.

To see additive inverses: given $x \in W$, take $\alpha = -1$, $\beta = 0$. Then $(-1)x + 0 = -x \in W$.

This single-condition test is sometimes more convenient, especially when verifying that a solution set (e.g., $\{x : Ax = 0\}$) is a subspace.
}

\bigskip

\hfill \break  $~\!\!$ \dotfill \medskip \\
\textbf{Question.}\\
Suppose
\begin{align*}
x = a_1 x_1 + \cdots + a_n x_n = b_1 x_1 + \cdots + b_n x_n.
\end{align*}
Show that $a_i = b_i$ for all $i$ using linear independence.

\noindent\\
\textbf{Solution.}\\
\quad\\
\noindent\fbox{ \parbox{\textwidth}{
Subtract the two expressions:
\begin{align*}
0
&=(a_1x_1+\cdots+a_nx_n)-(b_1x_1+\cdots+b_nx_n) \\
&=(a_1-b_1)x_1+\cdots+(a_n-b_n)x_n.
\end{align*}
If $\{x_1,\dots,x_n\}$ is linearly independent, then the only way a linear combination of these vectors can equal $0$ is if every coefficient is $0$. Hence
\begin{align*}
a_1-b_1=0,\quad \dots,\quad a_n-b_n=0.
\end{align*}
Therefore
\begin{align*}
a_i=b_i\qquad \text{for all }i.
\end{align*}
}}

\medskip

\bluebox{
\textbf{Alternative solution.} Use the coordinate map. The map $\Phi_B : V \to F^n$ given by $\Phi_B(x) = [x]_B$ is well-defined precisely when each $x$ has a unique coordinate vector. By assumption, $x$ maps to both $(a_1, \dots, a_n)$ and $(b_1, \dots, b_n)$. Since $\Phi_B$ is injective (its kernel is trivial: $\Phi_B(x) = 0$ implies $x = 0$), the output must be unique. Therefore $(a_1, \dots, a_n) = (b_1, \dots, b_n)$.

In other words, uniqueness of coordinates is equivalent to injectivity of the coordinate map, which is equivalent to linear independence of the basis vectors.
}

\bigskip

\hfill \break  $~\!\!$ \dotfill \medskip \\
\textbf{Question.}\\
Prove that the coordinate map $x \mapsto [x]_B$ is linear, injective, and surjective.

\noindent\\
\textbf{Solution.}\\
\quad\\
\noindent\fbox{ \parbox{\textwidth}{
Let $B=\{x_1,\dots,x_n\}$. Define $\Phi(x)=[x]_B$.

To prove linearity, write
\begin{align*}
x=a_1x_1+\cdots+a_nx_n,\qquad y=b_1x_1+\cdots+b_nx_n.
\end{align*}
Then
\begin{align*}
\Phi(x+y)=[x+y]_B=(a_1+b_1,\dots,a_n+b_n)=\Phi(x)+\Phi(y),
\end{align*}
and for $c\in F$,
\begin{align*}
\Phi(cx)=[cx]_B=(ca_1,\dots,ca_n)=c\Phi(x).
\end{align*}
So $\Phi$ is linear.

To prove injectivity, suppose $\Phi(x)=\Phi(y)$. Then $[x]_B=[y]_B$, so $x$ and $y$ have the same coordinates with respect to $B$. By uniqueness of basis expansions, $x=y$.

To prove surjectivity, let $(a_1,\dots,a_n)\in F^n$. Set
\begin{align*}
x=a_1x_1+\cdots+a_nx_n.
\end{align*}
Then
\begin{align*}
\Phi(x)=[x]_B=(a_1,\dots,a_n).
\end{align*}
Hence every vector in $F^n$ is hit by $\Phi$.

Therefore $\Phi$ is linear, injective, and surjective.
}}

\medskip

\bluebox{
\textbf{Alternative solution.} After establishing linearity, one can use the rank--nullity theorem to avoid checking surjectivity separately.

\textbf{Injectivity via kernel:} $\Phi(x) = 0$ means $(a_1, \dots, a_n) = (0, \dots, 0)$, so $x = 0$. Hence $\ker \Phi = \{0\}$.

\textbf{Surjectivity via dimension:} By rank--nullity, $\dim(\operatorname{Im} \Phi) = \dim V - \dim \ker \Phi = n - 0 = n$. Since $\operatorname{Im} \Phi \subseteq F^n$ and both have dimension $n$, we conclude $\operatorname{Im} \Phi = F^n$.

This is the standard shortcut: for a linear map between spaces of equal finite dimension, injectivity and surjectivity are equivalent.
}

\bigskip


\hfill \break  $~\!\!$ \dotfill \medskip \\
\textbf{Question.}\\
Construct the matrix $Q$ explicitly by expressing each vector of $B'$ in terms of $B$. Show that $Q$ is invertible.

\noindent\\
\textbf{Solution.}\\
\quad\\
\noindent\fbox{ \parbox{\textwidth}{
Let
\begin{align*}
B=\{x_1,\dots,x_n\},\qquad B'=\{x_1',\dots,x_n'\}.
\end{align*}
Because $B$ is a basis, each $x_j'$ can be written uniquely as
\begin{align*}
x_j'=a_{1j}x_1+\cdots+a_{nj}x_n.
\end{align*}
Define the matrix
\begin{align*}
Q=\begin{bmatrix}
a_{11} & a_{12} & \cdots & a_{1n} \\
a_{21} & a_{22} & \cdots & a_{2n} \\
\vdots & \vdots & \ddots & \vdots \\
a_{n1} & a_{n2} & \cdots & a_{nn}
\end{bmatrix}.
\end{align*}
The $j$th column of $Q$ is exactly $[x_j']_B$.

Now let
\begin{align*}
x=c_1x_1'+\cdots+c_nx_n'.
\end{align*}
Substituting the expressions for the $x_j'$ in terms of $B$ shows that the coordinate vector $[x]_B$ is obtained from $[x]_{B'}$ by multiplication by $Q$:
\begin{align*}
[x]_B=Q[x]_{B'}.
\end{align*}
So $Q$ is the change of coordinate matrix from $B'$-coordinates to $B$-coordinates.

To show $Q$ is invertible, suppose $Qc=0$ for some $c=(c_1,\dots,c_n)^T$. Then
\begin{align*}
0=Qc=[c_1x_1'+\cdots+c_nx_n']_B.
\end{align*}
Hence
\begin{align*}
c_1x_1'+\cdots+c_nx_n'=0.
\end{align*}
Since $B'$ is linearly independent, all $c_j=0$. Thus $\ker(Q)=\{0\}$, so $Q$ is injective. Because $Q$ is an $n\times n$ matrix, injective implies invertible. Therefore $Q$ is invertible.
}}

\medskip

\bluebox{
\textbf{Alternative solution.} Use the composition of coordinate isomorphisms. Define $Q = \Phi_B \circ \Phi_{B'}^{-1} : F^n \to F^n$. This is a composition of two linear isomorphisms, hence itself a linear isomorphism. Representing it as a matrix gives an invertible matrix.

Concretely, to find the $j$th column of $Q$, apply $Q$ to $e_j$ (the $j$th standard basis vector):
\[
Q e_j = \Phi_B(\Phi_{B'}^{-1}(e_j)) = \Phi_B(x_j') = [x_j']_B.
\]
So the columns of $Q$ are $[x_1']_B, \dots, [x_n']_B$, matching the explicit construction. Invertibility is immediate from the fact that a composition of isomorphisms is an isomorphism.
}

\bigskip


\hfill \break  $~\!\!$ \dotfill \medskip \\
\textbf{Question.}\\
Show that the matrix representation $[T]_B$ satisfies
\begin{align*}
[T(x)]_B = [T]_B [x]_B
\end{align*}
for all $x \in V$.

\noindent\\
\textbf{Solution.}\\
\quad\\
\noindent\fbox{ \parbox{\textwidth}{
Let $B=\{x_1,\dots,x_n\}$ be a basis of $V$. By definition, the $j$th column of $[T]_B$ is the coordinate vector of $T(x_j)$ with respect to $B$:
\begin{align*}
[T(x_j)]_B=\text{column }j\text{ of }[T]_B.
\end{align*}

Now let
\begin{align*}
x=a_1x_1+\cdots+a_nx_n.
\end{align*}
Then
\begin{align*}
[x]_B=
\begin{bmatrix}
a_1\\
\vdots\\
a_n
\end{bmatrix}.
\end{align*}
Since $T$ is linear,
\begin{align*}
T(x)=a_1T(x_1)+\cdots+a_nT(x_n).
\end{align*}
Taking coordinates with respect to $B$ gives
\begin{align*}
[T(x)]_B
&=a_1[T(x_1)]_B+\cdots+a_n[T(x_n)]_B.
\end{align*}
But multiplying the matrix $[T]_B$ by the vector $[x]_B$ produces exactly this same linear combination of the columns of $[T]_B$. Hence
\begin{align*}
[T(x)]_B=[T]_B[x]_B.
\end{align*}
This is the defining relation between a linear map and its matrix representation with respect to the basis $B$.
}}

\medskip

\bluebox{
\textbf{Alternative solution (via the commutative diagram).} From the lecture notes, the matrix $[T]_B$ is defined to make the following square commute:
\[
\begin{array}{ccc}
V & \xrightarrow{T} & V \\
\downarrow [\cdot]_B & & \downarrow [\cdot]_B \\
F^n & \xrightarrow{[T]_B} & F^n
\end{array}
\]
Formally, $[T]_B = \Phi_B \circ T \circ \Phi_B^{-1}$, where $\Phi_B(x) = [x]_B$. Then for any $x \in V$:
\[
[T]_B [x]_B = (\Phi_B \circ T \circ \Phi_B^{-1})(\Phi_B(x)) = \Phi_B(T(x)) = [T(x)]_B.
\]
This one-line computation encapsulates the entire proof. The matrix $[T]_B$ is simply the coordinate representation of $T$: convert from coordinates to abstract vectors ($\Phi_B^{-1}$), apply $T$, then convert back to coordinates ($\Phi_B$).
}


\bigskip

\hfill \break  $~\!\!$ \dotfill \medskip \\
\textbf{Question.}\\
Let $X := (X, \|\cdot \|_X$) and $Y= (Y, \|\cdot \|_Y$) be normed spaces, and let $T: X\rightarrow Y$ be a linear operator. Prove that if $\dim(X)<\infty$, then $T$ is bounded.
\noindent\\
\textbf{Solution.}\\
\quad\\
\noindent\fbox{ \parbox{\textwidth}{
Suppose $\dim(X) = n < \infty$, and let $\{e_1, \dots, e_n\}$ be a basis for $X$. Then for any $x \in X$, there exist unique scalars $a_1, \dots, a_n$ such that
\[
x = \sum_{i=1}^n a_i e_i.
\]
By linearity of $T$, we have
\[
T(x) = \sum_{i=1}^n a_i T(e_i).
\]
Using the triangle inequality,
\[
\|T(x)\|_Y \le \sum_{i=1}^n |a_i| \, \|T(e_i)\|_Y.
\]
Let
\[
M := \max_{1 \le i \le n} \|T(e_i)\|_Y.
\]
Then
\[
\|T(x)\|_Y \le M \sum_{i=1}^n |a_i|.
\]

Define a new norm on $X$ based on the basis $\{e_1, \dots, e_n\}$ by
\[
\|x\|_* := \sum_{i=1}^n |a_i|.
\]
Since $X$ is finite-dimensional, all norms on $X$ are equivalent. Hence there exists a constant $C > 0$ such that
\[
\|x\|_* \le C \|x\|_X \quad \text{for all } x \in X.
\]
Therefore,
\[
\|T(x)\|_Y \le M \|x\|_* \le MC \|x\|_X.
\]
Let $K := MC$. Then for all $x \in X$,
\[
\|T(x)\|_Y \le K \|x\|_X.
\]

Thus $T$ is bounded.

\hfill $\square$

}}

\medskip

\bluebox{
\textbf{Alternative solution (via compactness).} Let $\{e_1, \dots, e_n\}$ be a basis for $X$. The unit sphere $S = \{x \in X : \|x\|_X = 1\}$ is compact in finite dimensions (by the Heine--Borel theorem, since $S$ is closed and bounded).

The function $f : X \to \mathbb{R}$ defined by $f(x) = \|T(x)\|_Y$ is continuous (since $T$ is linear, hence continuous in finite dimensions, and $\|\cdot\|_Y$ is continuous). By the extreme value theorem, $f$ attains its maximum on $S$:
\[
K := \max_{x \in S} \|T(x)\|_Y < \infty.
\]
For any $x \neq 0$, we have $x / \|x\|_X \in S$, so
\[
\|T(x)\|_Y = \|x\|_X \cdot \left\|T\!\left(\frac{x}{\|x\|_X}\right)\right\|_Y \le K \|x\|_X.
\]
This also holds for $x = 0$. Hence $T$ is bounded with $\|T\| \le K$.

Note: this argument uses the fact that in finite dimensions, all norms are equivalent and closed bounded sets are compact --- both of which fail in infinite dimensions.
}

\bigskip

\hfill \break  $~\!\!$ \dotfill \medskip \\
\textbf{Question.}\\
Let $T: X\rightarrow Y$ be a linear operator between normed linear spaces. Suppose that $T$ is continuous. Then, show that $T(0)=0$.
\noindent\\
\textbf{Solution.}\\
\quad\\
\noindent\fbox{ \parbox{\textwidth}{
Since $T$ is linear, for any $x \in X$ we have
\[
T(0) = T(x - x) = T(x) - T(x) = 0.
\]
Thus $T(0) = 0$.

\medskip

Alternatively, using continuity: since $T$ is continuous at $0$, for any sequence $\{x_n\}$ in $X$ with $x_n \to 0$, we have $T(x_n) \to T(0)$. Taking $x_n = 0$ for all $n$, we get $T(x_n) = 0$ for all $n$, hence $T(0) = 0$.

\hfill $\square$
}}

\medskip

\bluebox{
\textbf{Alternative solution (via homogeneity).} Since $T$ is linear, it is homogeneous: $T(\alpha x) = \alpha T(x)$ for all $\alpha \in F$, $x \in X$. In particular, taking $\alpha = 0$ for any $x \in X$:
\[
T(0) = T(0 \cdot x) = 0 \cdot T(x) = 0.
\]
Note that this argument uses only linearity, not continuity. In fact, $T(0) = 0$ holds for \emph{every} linear operator, continuous or not. The continuity hypothesis in the problem statement is stronger than needed for this conclusion. However, the result is important as a stepping stone: the proof that continuity implies boundedness (from the lecture notes) begins by noting $T(0) = 0$ and then exploiting continuity at $0$.
}

\bigskip

\hfill \break  $~\!\!$ \dotfill \medskip \\
\textbf{Question.}\\
Let $T: X \to Y$ be a linear operator between normed linear spaces. Suppose that $T$ is continuous. Then, for $\varepsilon = 1$, there exists $\delta > 0$ such that
\[
\|u - 0\| < \delta \implies \|T(u) - T(0)\| < 1.
\]
Since $T(0) = 0$, this simplifies to
\[
\|u\| < \delta \implies \|T(u)\| < 1.
\]

Show that this implies
\[
\|u\| < \delta \implies \|T(u)\| \le 1.
\]
\noindent\\
\textbf{Solution.}\\
\quad\\
\noindent\fbox{ \parbox{\textwidth}{
Let $u \in X$ with $\|u\| < \delta$. For any $0 < t < 1$, we have
\[
\|t u\| = t \|u\| < \|u\| < \delta.
\]
Hence, by the previous implication,
\[
\|T(tu)\| < 1.
\]
Using linearity of $T$, we obtain
\[
\|t T(u)\| < 1 \quad \Rightarrow \quad t \|T(u)\| < 1.
\]
Thus,
\[
\|T(u)\| < \frac{1}{t}.
\]
Since this holds for all $0 < t < 1$, letting $t \to 1^{-}$ gives
\[
\|T(u)\| \le 1.
\]

\hfill $\square$

}}

\medskip

\bluebox{
\textbf{Alternative solution (direct from the open condition).} The set $\{u : \|T(u)\| < 1\}$ is open (as the preimage of the open set $(-\infty, 1)$ under the continuous function $u \mapsto \|T(u)\|$). It contains the open ball $B(0, \delta)$.

Now take any $u$ with $\|u\| < \delta$. Consider the function $g(s) = \|T(su)\|= s\|T(u)\|$ for $s \ge 0$. We know $g(s) < 1$ for all $s \in [0,1)$ (since $\|su\| < \delta$). Since $g$ is continuous and $g(s) < 1$ for all $s < 1$, by continuity $g(1) = \|T(u)\| \le 1$.

Equivalently, this is the general fact that if $f$ is continuous and $f(x) < c$ on a set approaching a limit point, then $f \le c$ at the limit. The strict inequality becomes weak at the boundary.
}

\bigskip

\hfill \break  $~\!\!$ \dotfill \medskip \\
\textbf{Question.}\\
Let $T, S : X \to Y$ be linear operators. Show that $\alpha T + \beta S$ is a linear operator for any scalars $\alpha, \beta$, and that the collection of all linear operators from $X$ to $Y$ forms a vector space.


\vspace{0.5cm}
\noindent
\textbf{Solution.}\\
\quad\\
\noindent\fbox{ \parbox{\textwidth}{
Let $u,v \in X$ and let $c \in F$. Define $(\alpha T+\beta S):X\to Y$ by
\begin{align*}
(\alpha T+\beta S)(x)=\alpha T(x)+\beta S(x).
\end{align*}
This is the pointwise definition of addition and scalar multiplication of functions. We can do this because $T(x)$ and $S(x)$ are elements of $Y$ since $T,S:X\to Y$; $Y$ is a vector space over $F$, so $\alpha T(x)$ and $\beta S(x)$ are in $Y$; and since $Y$ is closed under addition, $\alpha T(x)+\beta S(x)\in Y$. \\

We check linearity:
\begin{align*}
(\alpha T+\beta S)(u+v)
&=\alpha T(u+v)+\beta S(u+v) \\
&=\alpha(T(u)+T(v))+\beta(S(u)+S(v)) \\
&=(\alpha T(u)+\beta S(u))+(\alpha T(v)+\beta S(v)) \\
&=(\alpha T+\beta S)(u)+(\alpha T+\beta S)(v),
\end{align*}
and
\begin{align*}
(\alpha T+\beta S)(cu)
&=\alpha T(cu)+\beta S(cu) \\
&=\alpha cT(u)+\beta cS(u) \\
&=c(\alpha T(u)+\beta S(u)) \\
&=c(\alpha T+\beta S)(u).
\end{align*}
Hence $\alpha T+\beta S$ is linear.

Now let $\mathcal{L}(X,Y)$ denote the set of all linear operators from $X$ to $Y$. We show it is a vector space over $F$ with the usual addition and scalar multiplication.

First, by the calculation above, $\mathcal{L}(X,Y)$ is closed under addition and scalar multiplication. Also the zero map $0:X\to Y$, defined by $0(x)=0$ for all $x\in X$, is linear, so $0\in \mathcal{L}(X,Y)$. If $T\in \mathcal{L}(X,Y)$, then $(-1)T=-T$ is linear, so additive inverses exist.

The remaining vector space axioms follow pointwise from the corresponding axioms in $Y$. For example, for $T,S,R\in \mathcal{L}(X,Y)$ and $x\in X$,
\begin{align*}
((T+S)+R)(x)=T(x)+S(x)+R(x)=(T+(S+R))(x),
\end{align*}
so $(T+S)+R=T+(S+R)$. Similarly $T+S=S+T$, and for $a,b\in F$,
\begin{align*}
(a+b)T=aT+bT,\qquad a(T+S)=aT+aS,\qquad (ab)T=a(bT),\qquad 1T=T.
\end{align*}
Therefore $\mathcal{L}(X,Y)$ is a vector space.
}}

\medskip

\bluebox{
\textbf{Alternative solution (subspace test).} Rather than verifying all vector space axioms from scratch, use the subspace test from the lecture notes. Let $Y^X$ denote the vector space of \emph{all} functions from $X$ to $Y$ (with pointwise operations). Then $\mathcal{L}(X,Y) \subseteq Y^X$.

We verify $\mathcal{L}(X,Y)$ is a subspace of $Y^X$:
\begin{enumerate}[leftmargin=1.5em]
\item \textbf{Nonempty:} The zero map $0 \in \mathcal{L}(X,Y)$.
\item \textbf{Closure under addition:} If $T, S$ are linear, then $(T+S)(u+v) = T(u+v) + S(u+v) = T(u)+T(v)+S(u)+S(v) = (T+S)(u) + (T+S)(v)$, and similarly for scalar multiplication. So $T+S$ is linear.
\item \textbf{Closure under scalar multiplication:} If $T$ is linear, then $(\alpha T)(u+v) = \alpha T(u+v) = \alpha(T(u)+T(v)) = \alpha T(u) + \alpha T(v) = (\alpha T)(u) + (\alpha T)(v)$. So $\alpha T$ is linear.
\end{enumerate}

By the subspace test, $\mathcal{L}(X,Y)$ is a subspace of $Y^X$, hence a vector space. This approach avoids re-checking associativity, commutativity, distributivity, etc., since they are inherited from $Y^X$.

In finite dimensions, $\dim(\mathcal{L}(X,Y)) = \dim(X) \cdot \dim(Y)$, since choosing a matrix representation identifies $\mathcal{L}(X,Y) \cong F^{m \times n}$.
}

\end{document}
